Em Mauss e Beuchat, p. 436: \enquote{com efeito, há um notável sistema de tabu sobre o nome dos mortos entre os Esquimós, e esse tabu se observa por assentamento; disso resulta a supressão radical de todos os nomes comuns contidos nos nomes próprios dos indivíduos [falecidos]. Há, a seguir, um costume regular, no assentamento, de dar o nome do último morto ao primeiro a nascer; a criança é vista como o falecido reencarnado, e assim cada localidade possui um número determinado de nomes próprios que constituem, portanto, um elemento de sua fisionomia.}
Em Santos, p. 102: \enquote{somos povos de trajetórias, não somos povos de teoria. Somos da circularidade: começo, meio e começo. As nossas vidas não têm fim. A geração avó é o começo, a geração mãe é o meio e a geração neta é o começo de novo.}